\subsubsection{Valuation rings}
\label{editorial-source-section-valuation-rings}

\noindent
Here are some definitions.

\begin{definition}
\label{editorial-source-definition-valuation-ring}
Valuation rings.
\begin{enumerate}
\item Let $K$ be a field. Let $A$, $B$ be local rings contained
in $K$. We say that $B$ {\it dominates} $A$ if $A \subset B$
and $\mathfrak m_A = A \cap \mathfrak m_B$.
\item Let $A$ be a ring. We say $A$ is a {\it valuation ring}
if $A$ is a local domain and if $A$ is maximal
for the relation of domination among local rings contained in
the fraction field of $A$.
\item Let $A$ be a valuation ring with fraction field $K$.
If $R \subset K$ is a subring of $K$, then we say $A$
is {\it centered} on $R$ if $R \subset A$.
\end{enumerate}
\end{definition}

\noindent
With this definition a field is a valuation ring.

\begin{lemma}
\label{editorial-source-lemma-dominate}
Let $K$ be a field. Let $A \subset K$ be a local subring.
Then there exists a valuation ring with fraction field $K$
dominating $A$.
\end{lemma}

\begin{proof}
Consider the set of local subrings
$B\subseteq K$ dominating the
original $A$, ordered by domination.
It is nonempty since it contains
$A$. The relation is transitive:
if $B$ dominates $A$ and $C$
dominates $B$, then
$A\cap\mathfrak m_C=A\cap\mathfrak m_B=\mathfrak m_A$.
Antisymmetry follows from the
inclusions of the actual subrings.

For a nonempty chain $(B_i)$,
put $B=\bigcup_i B_i$ and
$\mathfrak n=\bigcup_i\mathfrak m_{B_i}$.
The maximal ideals are compatible
under contraction and inclusion.
A common chain member containing
any finite list of elements proves
that $B$ is a ring and
$\mathfrak n$ is an ideal.
It is proper because no member
contains $1$. If $b\in B\setminus\mathfrak n$,
it is a unit in any member
containing it, hence in $B$.
If $b\in\mathfrak n$, an inverse
in $B$ would put both elements
in a common later member and
contradict properness of its
maximal ideal. Thus $B$ is
local with maximal ideal
$\mathfrak n$, and
$B_i\cap\mathfrak n=\mathfrak m_{B_i}$
for every $i$. It dominates
all members and the original
$A$. The empty chain has
upper bound $A$. Zorn's lemma
therefore gives a maximal
member $V$.

We prove that its actual
fraction field $F=\operatorname{Frac}(V)$
equals $K$. Otherwise choose
$t\in K\setminus F$.
If $t$ is transcendental over
$F$, the evaluation map embeds
the original polynomial ring
$V[T]$ in $K$, with $T\mapsto t$.
The ideal $(\mathfrak m_V,t)$
of $V[t]$ is maximal, since
its quotient is $V/\mathfrak m_V$.
Its localization embeds in
$K$ by the original fraction
map and dominates $V$:
the maximal ideal contracts
to $\mathfrak m_V$. It contains
$t\notin V$, contradicting
maximality of $V$.

If $t$ is algebraic over
$F$, retain its monic equation
$$
t^d+c_1t^{d-1}+\cdots+c_d=0,
\qquad c_i\in F.
$$
Choose a nonzero $a\in V$
that is a product of denominators
for all original $c_i$, and
write $b_i=ac_i\in V$.
The original element $u=at$
satisfies the full equation
$$
u^d+b_1u^{d-1}+a b_2u^{d-2}
+\cdots+a^{d-1}b_d=0.
$$
This is exactly $a^d$ times
the displayed original equation,
using $a^{i-1}b_i=a^ic_i$;
every coefficient is retained.
Hence $V[u]$ is finite integral
over $V$, generated as a module
by $1,u,\ldots,u^{d-1}$.
By Lemma
\ref{editorial-source-lemma-integral-overring-surjective},
there is a prime $\mathfrak q$
of $V[u]$ contracting to
$\mathfrak m_V$. The original
fraction map embeds $V[u]_{\mathfrak q}$
in $K$, with maximal-ideal
contraction $\mathfrak m_V$.
It dominates $V$ and contains
$u\notin V$, since $u\in V$
would imply $t=u/a\in F$.
This is again a contradiction.
Thus $F=K$.

Every local subring of $K$
dominating $V$ also dominates
$A$. Maximality in the chosen
set therefore makes $V$ maximal
for domination in its own
fraction field $K$, exactly
the definition of a valuation
ring. This proves the lemma.

In particular, for any original
subring $R\subseteq K$ and
prime $\mathfrak p\subset R$,
apply the result to the
actual local subring $R_{\mathfrak p}$
of $K$. The resulting $V$
has fraction field $K$ and
$$
\mathfrak m_V\cap R
=(\mathfrak m_V\cap R_{\mathfrak p})\cap R
=\mathfrak pR_{\mathfrak p}\cap R
=\mathfrak p.
$$
Thus the center can be
prescribed as any given
prime of the original subring.
\end{proof}

\begin{lemma}
\label{editorial-source-lemma-valuation-ring-normal}
Let $A$ be a valuation ring.
Then $A$ is a normal domain.
\end{lemma}

\begin{proof}
Let $K=\operatorname{Frac}(A)$, retaining
the original fraction embedding. Suppose
$x\in K$ is integral over $A$.
The original $A'=A[x]\subseteq K$
is finite integral over $A$: a
monic equation of degree $d$ for
$x$ expresses all its powers in
the span of $1,x,\ldots,x^{d-1}$.
Lemma \ref{editorial-source-lemma-integral-overring-surjective}
gives a prime $\mathfrak m'$ of
$A'$ contracting to $\mathfrak m_A$.
The actual fraction map
$A'_{\mathfrak m'}\to K$, $a/s\mapsto a/s$,
is injective because $A'$ is a
domain and $s\notin\mathfrak m'$
is nonzero. The maximal ideal
of this localization contracts to
$\mathfrak m'$ in $A'$ and then
to $\mathfrak m_A$ in $A$.
Thus it dominates the original
$A$. The valuation-ring maximality
in this same $K$ gives
$A=A'_{\mathfrak m'}$, so $x\in A$.
Every element of the fraction field
integral over $A$ therefore belongs
to $A$, proving normality.
In particular any integral overring
$A\subseteq D\subseteq K$ equals
$A$, by applying this proof to
each of its original elements.
\end{proof}


\begin{lemma}
\label{editorial-source-lemma-valuation-ring-x-or-x-inverse}
Let $A$ be a valuation ring with maximal ideal $\mathfrak m$ and
fraction field $K$.
Let $x \in K^*$. Then either $x \in A$ or $x^{-1} \in A$ or both.
\end{lemma}

\begin{proof}
It suffices to treat the case
$x\notin A$; then $x\neq0$,
since $0\in A$. Retain the
original subring $A'=A[x]\subseteq K$.
If a prime $\mathfrak q$ of
$A'$ contracted to $\mathfrak m$,
the original fraction map would
embed $A'_{\mathfrak q}$ in $K$
as a local ring dominating
$A$. Maximality of the valuation
ring would give $A=A'_{\mathfrak q}$,
contrary to $x\notin A$.
There is therefore no such
prime. Every prime containing
$\mathfrak m A'$ would contract
to the maximal ideal $\mathfrak m$,
so $\mathfrak m A'$ is the
unit ideal. A finite expression
for its unit, expanded in
the original generator $x$,
gives
$$
1=\sum_{i=0}^d t_i x^i,\qquad t_i\in\mathfrak m.
$$
Here $d\geq1$, since $d=0$
would imply $1=t_0\in\mathfrak m$.
With $y=x^{-1}$, multiplication
by the full $x^{-d}$ gives
$$
(1-t_0)y^d-\sum_{i=1}^d t_i y^{d-i}=0.
$$
The index starts at $1$:
the original $t_0$ contribution
is already in $1-t_0$.
That coefficient is a unit
of the local $A$, so the
monic polynomial
$$
Q(T)=T^d-\sum_{i=1}^d(1-t_0)^{-1}t_iT^{d-i}
$$
annihilates $y$. The complete
original polynomial is
$(1-t_0)Q(T)$; its unit
factor and all coefficients
are kept in the comparison.
Thus $y$ is integral over
$A$ and belongs to $A$ by
Lemma \ref{editorial-source-lemma-valuation-ring-normal}.
This proves the assertion.

The inverse is in fact in
$\mathfrak m$: if $y$ were
a unit of $A$, then
$x=y^{-1}$ would belong to
$A$, contrary to the case
under consideration. Conversely,
the inverse of a nonzero
element of $\mathfrak m$
cannot belong to $A$, since
their product would put $1$
in $\mathfrak m$. Therefore
$$
K\setminus A=\{a^{-1}:0\neq a\in\mathfrak m\},
$$
and a nonzero $x$ and its
inverse both belong to $A$
exactly when $x\in A^*$.
The zero element belongs to
$A$ independently and is
never assigned an inverse.
\end{proof}

\begin{lemma}
\label{editorial-source-lemma-x-or-x-inverse-valuation-ring}
Let $A \subset K$ be a subring of a field $K$ such that for all
$x \in K^*$ either $x \in A$ or $x^{-1} \in A$ or both.
Then $A$ is a valuation ring with fraction field $K$.
\end{lemma}

\begin{proof}
The original $A$ is a nonzero
domain, as a unital subring
of the field $K$. The given
alternative proves that its
fraction field is all of $K$:
for nonzero $x\in K$, either
$x=x/1$ with $x\in A$ or
$x=1/x^{-1}$ with $x^{-1}\in A$.
These expressions are the
inverse of the canonical
injection $\operatorname{Frac}(A)\to K$.
The zero element is $0/1$.

If $A=K$, it is a local
ring with zero maximal ideal.
Otherwise it is not a field,
since a subfield equals its
own fraction field. Choose
a maximal ideal $\mathfrak m$
of $A$. If another maximal
ideal $\mathfrak m'$ were
distinct, their incomparability
would give
$$
x\in\mathfrak m\setminus\mathfrak m',\qquad
y\in\mathfrak m'\setminus\mathfrak m.
$$
Both original elements are
nonzero. If $x/y\in A$,
then $x=y(x/y)\in\mathfrak m'$,
a contradiction. If $y/x\in A$,
then $y=x(y/x)\in\mathfrak m$,
also a contradiction. The
given alternative excludes
two maximal ideals. Thus
$A$ is local, including
the field case.

Let the original local subring
$A'\subseteq K$ dominate $A$.
If $x\in A'\setminus A$,
then $x\neq0$ and the
hypothesis gives $x^{-1}\in A$.
It is not a unit of $A$,
since otherwise $x\in A$.
Hence $x^{-1}\in\mathfrak m_A$
and domination puts it in
$\mathfrak m_{A'}$. Multiplying
by the original $x\in A'$
would put $1$ in that ideal.
This contradiction gives
$A'=A$, proving the valuation
ring assertion in the exact
fraction field $K$.

This criterion also proves
that all ideals of a valuation
ring are linearly ordered
by inclusion. For nonzero
$a,b\in A$, applying the
alternative to $a/b$ gives
either $a=bc$ or $b=ac$
with $c\in A$. The zero
cases are immediate divisibility
equalities. Thus principal
ideals are comparable. If
ideals $I,J$ satisfy
$I\not\subseteq J$, choose
$a\in I\setminus J$.
For every $b\in J$, the
alternative $a=bc$ is
impossible because it puts
$a$ in $J$. The other
alternative gives $b\in(a)\subseteq I$;
this includes $b=0$.
Consequently $J\subseteq I$.
In particular all original
prime ideals are in one
inclusion chain. Conversely,
comparability of the original
principal ideals in a nonzero
domain gives the fraction
alternative in its own
fraction field, and the proof
above makes it a valuation ring.
\end{proof}

\begin{lemma}
\label{editorial-source-lemma-colimit-valuation-rings}
\begin{slogan}
Valuation rings are stable under filtered direct limits
\end{slogan}
Let $I$ be a directed set. Let $(A_i, \varphi_{ij})$
be a system of valuation rings over $I$.
Then $A = \colim A_i$ is a valuation ring.
\end{lemma}

\begin{proof}
The directed set is nonempty,
and all transition maps are
the original unital ring maps;
they need not be injective
or local. The colimit $A$
is nonzero: an equality
$1=0$ in it would hold
at a later stage, impossible
in a valuation domain $A_j$.
For $a,b\in A$ with $ab=0$,
choose representatives $a_i,b_i$
at one stage. At a later
stage their product is zero.
That stage is a domain,
so one of the two images
is zero and the corresponding
original $a$ or $b$ is zero.
Thus $A$ is a domain.

For arbitrary $a,b\in A$,
choose such representatives.
If either representative is
zero, one divides the other.
Otherwise Lemma
\ref{editorial-source-lemma-valuation-ring-x-or-x-inverse}
in $\operatorname{Frac}(A_i)$
gives either $a_i=b_i c_i$
or $b_i=a_i c_i$ with
$c_i\in A_i$. The same
original equality holds in
$A$ under the structure map.
For a fraction $a/b$ with
$b\neq0$, this proves that
$a/b\in A$ or, if $a\neq0$,
its inverse belongs to $A$.
Lemma
\ref{editorial-source-lemma-x-or-x-inverse-valuation-ring}
now proves the assertion.

There is an exact fraction-field
description even for the original
noninjective system. Put
$I_i=\ker(A_i\to A)$ and
$B_i=A_i/I_i$. Since $B_i$
injects in the domain $A$,
$I_i$ is prime and $B_i$
is a nonzero domain. The
divisibility equalities in
$A_i$ descend to $B_i$,
so the same criterion makes
$B_i$ a valuation ring.
The induced maps $B_i\to B_j$
are injective: their kernels
are zero because an element
vanishing in $B_j$ has zero
image in $A$, hence was
already in $I_i$.
These are the actual image
subrings of $A$, with
$A=\bigcup_i B_i$.

The field maps
$\operatorname{Frac}(B_i)\to\operatorname{Frac}(B_j)$
send each original $a_i/b_i$
to the ratio of its images;
denominators remain nonzero
by injectivity. They give
the isomorphism
$$
\mathop{\rm colim}_i\operatorname{Frac}(B_i)
\longrightarrow\operatorname{Frac}(A),
\qquad[i,a_i/b_i]\longmapsto a_i/b_i.
$$
Every fraction on the right
has both coefficients in
one common $B_i$, proving
surjectivity. Equality of two
fractions is the equality
of their cross products;
it holds in a common stage
because those stages inject
in $A$. This proves injectivity
and gives the inverse by
choosing a stage for the
two original coefficients.
No fraction-field map is
asserted for an original
transition map before its
kernel has been accounted for.

For a representative $a_i$,
its image in $A$ is a unit
exactly when its image is
a unit at some later stage.
One direction preserves the
original inverse equation.
For the other, an inverse
in $A$ has a representative
at a common stage, and
the equation giving product
$1$ holds at a later stage.
Hence the exact maximal-ideal
criterion is that every
later image of $a_i$ belong
to the corresponding maximal
ideal. A union of the
original maximal ideals need
not work: for the two-stage
system $V\to\operatorname{Frac}(V)$
with a nonfield valuation
ring $V$, each nonzero
element of $\mathfrak m_V$
becomes a unit at the
second stage. The colimit
is the field, with zero
maximal ideal, as the exact
criterion predicts.
\end{proof}

\begin{lemma}
\label{editorial-source-lemma-valuation-ring-cap-field}
Let $L/K$ be an extension of fields. If $B \subset L$
is a valuation ring, then $A = K \cap B$ is a valuation ring.
\end{lemma}

\begin{proof}
Retain the original fields $K\subseteq L$
and the original $B\subseteq L$.
Put $F=\operatorname{Frac}(B)\subseteq L$
and $E=K\cap F$. The ring in
the statement is exactly
$A=K\cap B=E\cap B$; these
are equal subrings of the
original $L$. For every
nonzero $x\in E$, the
valuation alternative in $B$
gives $x\in B$ or $x^{-1}\in B$.
Both $x,x^{-1}$ lie in $E$,
so the alternative holds in
$A$. Lemma
\ref{editorial-source-lemma-x-or-x-inverse-valuation-ring}
makes $A$ a valuation ring
whose exact fraction field
is $E=K\cap F$.
The canonical fraction map
$\operatorname{Frac}(A)\to E$
sends $a/b$ to the same
fraction in $L$. Its inverse
sends a nonzero $x$ to
$x/1$ if $x\in A$, or
to $1/x^{-1}$ otherwise;
the two possible expressions
agree when both apply.
This proves the fraction-field
identification without replacing
either original ambient field.

An original nonzero $a\in A$
is a unit of $A$ precisely
when it is a unit of $B$.
Indeed, if $a^{-1}\in B$,
then $a^{-1}\in K$ as well,
so it belongs to $A$.
The converse follows from
the inclusion. Hence
$$
\mathfrak m_A=A\cap\mathfrak m_B.
$$
The resulting residue map
$A/\mathfrak m_A\to B/\mathfrak m_B$
sends $[a]$ to $[a]$ and
is injective by this exact
kernel calculation. In particular
$B$ dominates the original
intersection $A$.

This also proves extension
of any given valuation ring
$A$ of a field $K$ to
any prescribed extension $L/K$.
Lemma \ref{editorial-source-lemma-dominate}
provides a valuation ring
$B\subseteq L$ of fraction
field $L$ dominating $A$.
Its intersection $A'=K\cap B$
has fraction field $K$ by
the proved formula and is
local with maximal ideal
$A'\cap\mathfrak m_B$.
It contains $A$ and its
maximal ideal contracts to
$\mathfrak m_A$, so it
dominates $A$. Maximality
of the original valuation
ring $A$ in $K$ forces
$A'=A$. Thus the extension
restricts to the original
ring and residue embedding
exactly, not merely to some
valuation ring containing it.
\end{proof}

\begin{lemma}
\label{editorial-source-lemma-valuation-ring-cap-field-finite}
Let $L/K$ be an algebraic extension of fields. If $B \subset L$
is a valuation ring with fraction field $L$ and not a field, then
$A = K \cap B$ is a valuation ring and not a field.
\end{lemma}

\begin{proof}
By Lemma \ref{editorial-source-lemma-valuation-ring-cap-field},
the original $A=K\cap B$ is
a valuation ring, and its
fraction field is exactly $K$
because $\operatorname{Frac}(B)=L$.
If $A$ were a field, it
would equal its fraction field,
so $A=K\subseteq B$.
Every $b\in B$ is algebraic
over $K$, hence integral over
this field by its monic
minimal polynomial. Thus
$K\subseteq B$ is integral.
Lemma \ref{editorial-source-lemma-integral-no-inclusion}
forbids a proper inclusion
of primes in $B$ over
the sole prime of $K$.
But the local domain $B$
has $(0)\subsetneq\mathfrak m_B$
because it is not a field,
a contradiction.

Equivalently the contradiction
has a direct inverse formula.
For nonzero $b\in B$, its
original monic minimal polynomial
$$
b^d+c_1b^{d-1}+\cdots+c_d=0
$$
has $c_d\neq0$ in $K$,
and yields
$$
b^{-1}=-c_d^{-1}
\left(b^{d-1}+\sum_{i=1}^{d-1}c_i b^{d-1-i}\right)\in B.
$$
For $d=1$ the parenthesis
is $1$. Thus every nonzero
$b$ would be a unit,
again making $B$ a field.

The algebraicity hypothesis
cannot be dropped. For any
field $k$, retain
$K=k$, $L=k(t)$ and
$B=k[t]_{(t)}$. Its fraction
field is the original $L$.
For a nonzero original
fraction $P(t)/Q(t)$, write
$P=t^rP_0$, $Q=t^sQ_0$
with $P_0(0),Q_0(0)\neq0$,
retaining both powers and
both residual factors.
If $r\geq s$, the fraction
$t^{r-s}P_0/Q_0$ lies in
$B$; if $r<s$, its inverse
$t^{s-r}Q_0/P_0$ lies in
$B$. The fraction criterion
therefore makes $B$ a
valuation ring. It is not
a field: evaluation at $t=0$
on this localization has
proper kernel $(t)B$, which
contains its nonzero $t$.
Nevertheless $K\cap B=k$
is a field. This identifies
the exact failure without
the original algebraicity
assumption.
\end{proof}

\begin{lemma}
\label{editorial-source-lemma-make-valuation-rings}
Let $A$ be a valuation ring. For any prime ideal $\mathfrak p \subset A$ the
quotient $A/\mathfrak p$ is a valuation ring. The same is true for the
localization $A_\mathfrak p$ and in fact any nonzero localization of $A$.
\end{lemma}

\begin{proof}
For the original prime
$\mathfrak p\subset A$, the
quotient $A/\mathfrak p$ is
a nonzero domain. Any two
nonzero residue classes have
original representatives $a,b$
outside $\mathfrak p$.
The valuation alternative in
$\operatorname{Frac}(A)$ gives
$a=bc$ or $b=ac$ for
some original $c\in A$.
The same equality in the
quotient gives its fraction
alternative, so Lemma
\ref{editorial-source-lemma-x-or-x-inverse-valuation-ring}
makes $A/\mathfrak p$ a
valuation ring. Its maximal
ideal is $\mathfrak m_A/\mathfrak p$:
a class outside that ideal
has an inverse inherited
from $A$. Its fraction field
is the original residue
field $\kappa(\mathfrak p)$,
through the inverse maps
$\overline a/\overline b\mapsto[a/b]$
and $[a/b]\mapsto\overline a/\overline b$
between $\operatorname{Frac}(A/\mathfrak p)$
and $A_{\mathfrak p}/\mathfrak pA_{\mathfrak p}$,
where $b\notin\mathfrak p$.

Let $W\subseteq A$ be an
original multiplicative subset.
If $0\in W$, Lemma
\ref{algebra-lemma-localization-zero}
gives $W^{-1}A=0$, which
is not a domain and hence
not a valuation ring.
This is an allowed localization
under Definition
\ref{algebra-definition-multiplicative-subset}.
If $0\notin W$, the
fraction map embeds $W^{-1}A$
in $K=\operatorname{Frac}(A)$:
equality of fractions is
exactly the original cross-product
equality, since nonzero elements
of $W$ do not annihilate
anything in the domain.
It contains $A$ and has
fraction field $K$. The
alternative $x\in A$ or
$x^{-1}\in A$ for every
nonzero $x\in K$ therefore
holds in $W^{-1}A$, making
it a valuation ring. This
includes $A_{\mathfrak p}$,
whose denominator set avoids zero.

More generally every original
overring $A\subseteq B\subseteq K$
is exactly $A_{\mathfrak q}$
for one prime $\mathfrak q$
of $A$. The same fraction
alternative first proves that
$B$ is a valuation ring
with fraction field $K$.
Put $\mathfrak q=A\cap\mathfrak m_B$.
The original fraction map
$A_{\mathfrak q}\to B$ is
defined because every denominator
outside $\mathfrak q$ is
a unit of $B$, and is
injective inside $K$.
For $b\in B$, if $b\in A$
it is already in $A_{\mathfrak q}$.
Otherwise $b\neq0$ and
$b^{-1}\in A$ by the
original valuation criterion.
It is a unit of $B$,
so $b^{-1}\notin\mathfrak q$
and $b=1/b^{-1}\in A_{\mathfrak q}$.
This proves equality of the
original subrings. The inverse
prime assignment is contraction
of the maximal ideal, since
$\mathfrak qA_{\mathfrak q}\cap A=\mathfrak q$.
Thus primes and overrings
correspond bijectively, reversing
inclusion: larger primes invert
fewer original elements.
Conversely an inclusion of overrings
preserves their units, so contraction
of their maximal ideals reverses
that inclusion as well.

For the given $W$, this
prime is explicitly
$$
\mathfrak q=\{a\in A:(a)\cap W=\emptyset\}.
$$
Indeed $a/1$ is a unit
in $W^{-1}A$ precisely when
$ab=w$ for some $b\in A$
and $w\in W$. An inverse
$b/w$ gives this exact
equality by cancelling the
nonzero localization witness
in the domain, and the
equality conversely gives
the inverse. This proves
the formula from the maximal
ideal of the valuation ring.
The prime is disjoint from
$W$ and contains every prime
disjoint from $W$: if
$a$ is in such a prime,
no multiple $ab$ can be
in $W$. It is therefore
the largest such prime,
and the canonical fraction
map identifies $W^{-1}A$
with exactly $A_{\mathfrak q}$.
\end{proof}

\begin{lemma}
\label{editorial-source-lemma-stack-valuation-rings}
Let $A'$ be a valuation ring with residue field $K$.
Let $A$ be a valuation ring with fraction field $K$.
Then
$C = \{\lambda \in A' \mid \lambda \bmod \mathfrak m_{A'} \in A\}$
is a valuation ring.
\end{lemma}

\begin{proof}
Keep the original residue map
$\pi:A'\to K$ and write
$\mathfrak p=\mathfrak m_{A'}$.
The original $C=\pi^{-1}(A)$
is a subring of $A'$,
and $\mathfrak p\subseteq C$.
The map $C/\mathfrak p\to A$
sending $[c]$ to $\pi(c)$
is an isomorphism: surjectivity
follows by lifting each $a\in A$
through $\pi$, and its kernel
is exactly $\mathfrak p$.
The inverse sends $a$ to
the class of any such lift;
different lifts differ in
$\mathfrak p$. In particular
$\mathfrak p$ is a prime
of the original $C$.

There is a more precise
comparison than just equality
of fraction fields:
$$
C_{\mathfrak p}\longrightarrow A',\qquad c/t\longmapsto c t^{-1}
$$
is an isomorphism inside
$F=\operatorname{Frac}(A')$.
It is defined and injective
because $t\in C\setminus\mathfrak p$
is a nonzero unit of $A'$.
For surjectivity, take the
original $b\in A'$ and
write its residue in
$K=\operatorname{Frac}(A)$
as $\pi(b)=a/d$ with
$a,d\in A$ and $d\neq0$;
for zero residue take $a=0,d=1$.
Choose $t\in A'$ lifting
this original $d$. Then
$t\in C\setminus\mathfrak p$
and $tb\in C$ because
$\pi(tb)=a$. Thus the
inverse sends $b$ to
$(tb)/t$. Independence of
the lift and denominator
follows from the equality
of cross products in the
original field $F$. Both
composites are identities.
It follows that
$\operatorname{Frac}(C)=F$,
including the case $A'$ is
a field and $\mathfrak p=0$.

For nonzero $x\in F$, if
$x\notin A'$, Lemma
\ref{editorial-source-lemma-valuation-ring-x-or-x-inverse}
gives $x^{-1}\in\mathfrak m_{A'}\subseteq C$.
If $x\in\mathfrak m_{A'}$,
then $x\in C$. In the
remaining case $x$ is a
unit of $A'$, so its
nonzero residue satisfies
$\pi(x)\in A$ or
$\pi(x)^{-1}\in A$ by the
valuation alternative in the
original residue field $K$.
Because $\pi(x^{-1})=\pi(x)^{-1}$,
this says $x\in C$ or
$x^{-1}\in C$. The criterion
of Lemma
\ref{editorial-source-lemma-x-or-x-inverse-valuation-ring}
therefore proves the assertion.

The maximal ideal and residue
field are also exact:
$$
\mathfrak m_C=\pi^{-1}(\mathfrak m_A),\qquad
C/\mathfrak m_C\simeq A/\mathfrak m_A.
$$
Indeed an element $c$ outside
that inverse image has
$\pi(c)\in A^*$, so it
is a unit of $A'$ and
$\pi(c^{-1})=\pi(c)^{-1}\in A$,
whence $c^{-1}\in C$.
An element inside cannot be
a unit, by applying $\pi$
to an inverse equation.
The residue isomorphism is
the map $[c]\mapsto[\pi(c)]$,
with inverse given by a
lift as above.

Finally the original prime
chain of $C$ is described
by these two exact maps.
Primes contained in $\mathfrak p$
correspond to primes of
$C_{\mathfrak p}=A'$ by
extension and contraction.
Primes containing $\mathfrak p$
correspond to primes of
$C/\mathfrak p=A$ by
quotient and inverse image.
All primes of the valuation
ring $C$ are comparable,
as proved in Lemma
\ref{editorial-source-lemma-x-or-x-inverse-valuation-ring},
so these two portions cover
its spectrum. They meet
at the single original prime
$\mathfrak p$, corresponding
to $\mathfrak m_{A'}$ on
the localization side and
to $(0)$ on the quotient
side. This proves the full
prime-chain composition.
\end{proof}

\begin{lemma}
\label{editorial-source-lemma-find-valuation-rings}
Let $A$ be a normal domain with fraction field $K$.
\begin{enumerate}
\item For every $x \in K$, $x \not \in A$ there exists a valuation ring
$A \subset V \subset K$ with fraction field $K$ such that $x \not \in V$.
\item If $A$ is local, we can moreover choose $V$ which dominates $A$.
\end{enumerate}
In other words, $A$ is the intersection of all valuation rings in $K$
containing $A$ and if $A$ is local, then $A$ is the intersection of
all valuation rings in $K$ dominating $A$.
\end{lemma}

\begin{proof}
For the given normal domain,
$x\in K\setminus A$ is
not integral over $A$.
We prove the construction
using exactly this nonintegrality
condition; it will also give
the integral-closure statement
for arbitrary domains below.
Since $0\in A$, the original
$x$ is nonzero. Put
$B=A[x^{-1}]\subseteq K$.
If $x\in B$, a finite
expression
$x=\sum_{i=0}^d a_i x^{-i}$
would give the full monic
equation
$$
x^{d+1}-\sum_{i=0}^d a_i x^{d-i}=0,
$$
contradicting nonintegrality.
Thus $x^{-1}$ is not a
unit of $B$, and
$x^{-1}B$ is a proper
ideal. Choose a prime
$\mathfrak p$ containing it
and apply Lemma
\ref{editorial-source-lemma-dominate} to the
original local subring
$B_{\mathfrak p}\subseteq K$.
It gives a valuation ring
$V$ of fraction field $K$
dominating $B_{\mathfrak p}$.
The element $x^{-1}$ belongs
to $\mathfrak m_V$, so
$x\notin V$: otherwise
their product would put
$1$ in $\mathfrak m_V$.
Since $A\subseteq B\subseteq V$,
this proves (1).

Suppose now the original
$A$ is local. Retain
$$
J=x^{-1}B+\mathfrak m_A B.
$$
If $J=B$, finite expressions
in the original $x^{-1}$,
with zero padding to a
common degree $d\geq0$,
give
$$
1=\left(\sum_{i=0}^d a_i x^{-i}\right)x^{-1}
+\sum_{i=0}^d a'_i x^{-i},
\qquad a_i\in A,\quad a'_i\in\mathfrak m_A.
$$
Multiplying by the full
$x^{d+1}$ and keeping the
two original coefficient lists
gives
$$
(1-a'_0)x^{d+1}
-\sum_{i=0}^d a_i x^{d-i}
-\sum_{i=1}^d a'_i x^{d+1-i}=0.
$$
For $d=0$ the last sum
is empty and the equation
is $(1-a'_0)x-a_0=0$.
In every case $u=1-a'_0$
is a unit of $A$.
The monic polynomial
$$
Q(T)=T^{d+1}
-\sum_{i=0}^d u^{-1}a_i T^{d-i}
-\sum_{i=1}^d u^{-1}a'_i T^{d+1-i}
$$
annihilates $x$, and the
complete polynomial originally
obtained is exactly $uQ(T)$.
Thus $x$ would be integral,
a contradiction. It follows
that the original $J$ is
proper. Choose a prime
$\mathfrak p\supseteq J$ and
then a valuation ring $V$
of fraction field $K$
dominating $B_{\mathfrak p}$,
as in the first construction.

The contraction $\mathfrak p\cap A$
contains $\mathfrak m_A$ and
is proper, so it equals
$\mathfrak m_A$. Localization
gives
$\mathfrak m_{B_{\mathfrak p}}\cap A=\mathfrak m_A$,
and domination gives
$\mathfrak m_V\cap A=\mathfrak m_A$.
Hence this original $V$
dominates $A$ and still
excludes $x$ because
$x^{-1}\in\mathfrak m_V$.
This proves (2). Each
asserted intersection equals
$A$: inclusion of $A$
is built into all its
members, and the construction
excludes every original element
of $K\setminus A$ from
at least one member.

More generally let $A$
be any domain with original
fraction field $K$, and
let $\widetilde A$ be its
integral closure in $K$.
Every $x\in\widetilde A$
belongs to every valuation
ring $A\subseteq V\subseteq K$:
its original monic equation
over $A$ is the same
monic equation over $V$,
and $V$ is normal by
Lemma \ref{editorial-source-lemma-valuation-ring-normal}.
If $x\notin\widetilde A$,
the construction above applies
because $x$ is not integral
over $A$, and gives a
member excluding it. Therefore
$$
\widetilde A=\bigcap_{A\subseteq V\subseteq K\text{ valuation}}V.
$$
For local $A$, exactly the
second construction gives
the same equality with the
intersection restricted to
valuation rings dominating $A$.

For a prescribed prime
$\mathfrak p\subset A$, the
valuation rings containing
$A$ with center
$\mathfrak m_V\cap A=\mathfrak p$
are precisely those dominating
the original $A_{\mathfrak p}$.
Indeed every $a\in A\setminus\mathfrak p$
is a unit of $V$, giving
the actual inclusion of
fractions $A_{\mathfrak p}\to V$,
and its maximal ideal
contracts to $\mathfrak pA_{\mathfrak p}$.
The converse follows by
contracting once more to
$A$. Their intersection
is thus exactly the
integral closure of
$A_{\mathfrak p}$ in $K$,
by the proved local formula.
\end{proof}

\noindent
A {\it totally ordered abelian group} is a pair $(\Gamma, \geq)$
consisting of an abelian group $\Gamma$ endowed with a total
ordering $\geq$ such that $\gamma \geq \gamma' \Rightarrow
\gamma + \gamma'' \geq \gamma' + \gamma''$ for all
$\gamma, \gamma', \gamma'' \in \Gamma$.

\begin{lemma}
\label{editorial-source-lemma-valuation-group}
Let $A$ be a valuation ring with field of fractions $K$.
Set $\Gamma = K^*/A^*$ (with group law written additively).
For $\gamma, \gamma' \in \Gamma$
define $\gamma \geq \gamma'$ if and only if
$\gamma - \gamma'$ is in the image of $A - \{0\} \to \Gamma$.
Then $(\Gamma, \geq)$ is a totally ordered abelian group.
\end{lemma}

\begin{proof}
Retain the original quotient
$\Gamma=K^*/A^*$ and write
$v(x)=[x]$ additively, so
$v(xy)=v(x)+v(y)$ by the
quotient group law. For original
$x,y\in K^*$, the proposed
order is exactly
$$
v(x)\geq v(y)\quad\Longleftrightarrow\quad x/y\in A.
$$
Indeed $v(x/y)$ is in the
image of $A\setminus\{0\}$
precisely when $x/y$ differs
from such an element by
a unit of $A$. Multiplication
by that unit and its inverse
proves the equivalence.
If the representatives are
changed to $xu,yw$ with
$u,w\in A^*$, the original
ratio becomes $(x/y)(u/w)$.
Its membership in $A$ is
unchanged, so the order is
well-defined on the quotient.

Reflexivity follows from $x/x=1$.
If $x/y,y/x\in A$, their
product is $1$, so $x/y$
is a unit and $v(x)=v(y)$;
this is antisymmetry. If
$x/y$ and $y/z$ belong
to $A$, their exact product
$x/z$ does too, proving
transitivity. For every pair,
the valuation alternative for
$x/y$ gives $x/y\in A$
or $y/x\in A$, proving
totality. Finally multiplying
both original representatives
by the same $z\in K^*$
does not change their ratio:
$(xz)/(yz)=x/y$. This
proves translation invariance
of the order and the
claimed ordered group structure.

Every totally ordered abelian
group is torsion-free. If
$\gamma>0$, repeated strict
translation invariance gives
$n\gamma>0$ for every integer
$n\geq1$; if $\gamma<0$,
apply the same argument to
$-\gamma$. Hence $n\gamma=0$
with $n\geq1$ forces
$\gamma=0$. In particular
this holds for the original
$\Gamma$.

The order recovers the
original ring and maximal
ideal exactly:
$$
A=\{0\}\cup\{x\in K^*:v(x)\geq0\},\qquad
\mathfrak m_A=\{0\}\cup\{x\in K^*:v(x)>0\}.
$$
The first equality is the
ratio criterion with $y=1$.
For $x\in A\setminus\{0\}$,
the quotient definition says
$v(x)=0$ exactly when
$x\in A^*$, giving the
second equality by localness.
Also $v(a/b)=v(a)-v(b)$
for all original nonzero
$a,b\in K$. The value
of $0$ is not introduced
into this group. If $A=K$,
the quotient is the zero
group, with its unique order.
Conversely, if this quotient
is zero, every nonzero
element of $K$ is a unit
of $A$, so $A=K$.
\end{proof}

\begin{definition}
\label{editorial-source-definition-value-group}
Let $A$ be a valuation ring.
\begin{enumerate}
\item The totally ordered abelian group $(\Gamma, \geq)$ of
Lemma \ref{editorial-source-lemma-valuation-group} is called the
{\it value group} of the valuation ring $A$.
\item The map $v : A - \{0\} \to \Gamma$ and also $v : K^* \to \Gamma$ is
called the {\it valuation} associated to $A$.
\item The valuation ring $A$ is called a {\it discrete valuation ring}
if $\Gamma \cong \mathbf{Z}$.
\end{enumerate}
\end{definition}

\noindent
Suppose the original ordered group
$\Gamma$ is isomorphic as a group
to $\mathbf Z$, and choose an
initial group isomorphism
$\psi:\Gamma\to\mathbf Z$.
Put $\gamma=\psi^{-1}(1)\neq0$.
Let $\varepsilon=1$ if
$\gamma>0$, and $\varepsilon=-1$
if $\gamma<0$. Then
$$
\phi=\varepsilon\psi:\Gamma\longrightarrow\mathbf Z
$$
is an order isomorphism with
the usual order on the
receiving integers. Indeed
$\varepsilon\gamma>0$, and
every original element is
$n\gamma$ for a unique
$n\in\mathbf Z$. Repeated
addition and translation
invariance give
$n\gamma\geq0$ precisely when
$\varepsilon n\geq0$, which
is precisely $\phi(n\gamma)\geq0$.
The same applies to differences,
proving both order implications.
Any other group isomorphism
is $\delta\psi$ for
$\delta\in\{1,-1\}$, since
the only generators of
$\mathbf Z$ are $1,-1$.
Requiring the inverse image
of $1$ to be positive
forces $\delta=\varepsilon$.
This is the precise unique
order-compatible isomorphism.
The original $\Gamma$ and
its value map $v$ are
retained; their integer comparison
is $\phi\circ v$, with
$v=\phi^{-1}\circ(\phi\circ v)$.

Write $\gamma_0=\phi^{-1}(1)>0$.
Surjectivity of the original
quotient $v:K^*\to\Gamma$
gives an element $\pi\in K^*$
with $v(\pi)=\gamma_0$.
The positive-cone description
puts $\pi$ in the maximal
ideal of $A$. For every
original nonzero $a\in A$,
put $n=\phi(v(a))\geq0$.
The exact original element
$u=a/\pi^n$ has value
$v(a)-n\gamma_0=0$, so
$u\in A^*$ and
$$
a=u\pi^n.
$$
Both the unit $u$ and
the original exponent are
part of this equality;
neither is discarded. For
the chosen $\pi$, the
integer $n$ is unique by
applying $\phi v$, and
then $u=a/\pi^n$ is
unique as well. Another
element $\pi'$ of the
same least positive value
has $\pi'/\pi\in A^*$,
which is the exact comparison
between such choices.

Every nonzero ideal $I$
has the form $(\pi^n)$
for exactly one $n\geq0$.
Choose the least integer
among $\phi(v(a))$ for
the original nonzero $a\in I$,
and choose such an $a$.
The full equality $a=u\pi^n$
with $u$ a unit puts
$\pi^n=u^{-1}a$ in $I$.
For every nonzero $b\in I$,
the inequality $v(b)\geq n\gamma_0$
puts $b/\pi^n$ in $A$,
and a zero member contributes
the equality $0=\pi^n0$.
Thus all of $I$ is
contained in $(\pi^n)$,
giving equality. Distinct
$n$ give distinct ideals
by their least original
values. The zero ideal
is retained as a separate
case and is not assigned
an infinite value in $\Gamma$.
In particular the maximal
ideal is $(\pi)$, every
ideal is principal, and
the only prime ideals are
$(0)$ and $(\pi)$: for
$n>1$, the product
$\pi\pi^{n-1}$ lies in
$(\pi^n)$ although neither
factor does. This proves
the ideal and prime descriptions
from the original ordered
quotient without suppressing
any unit factors.

\begin{lemma}
\label{editorial-source-lemma-properties-valuation}
Let $A$ be a valuation ring. The valuation $v : A -\{0\} \to \Gamma_{\geq 0}$
has the following properties:
\begin{enumerate}
\item $v(a) = 0 \Leftrightarrow a \in A^*$,
\item $v(ab) = v(a) + v(b)$,
\item $v(a + b) \geq \min(v(a), v(b))$ provided $a + b \not = 0$.
\end{enumerate}
\end{lemma}

\begin{proof}
The quotient map to the
original $\Gamma=K^*/A^*$
has kernel $A^*$, giving
(1), and its group law
gives (2). For (3), retain
the original nonzero $a,b$
and assume $a+b\neq0$.
Interchanging their roles
if necessary, suppose
$v(a)\leq v(b)$. The
order proved in Lemma
\ref{editorial-source-lemma-valuation-group}
then gives $b/a\in A$.
The exact equality
$$
\frac{a+b}{a}=1+\frac ba
$$
puts the nonzero left-hand
element in $A$, so
$v(a+b)-v(a)\geq0$.
This proves (3), retaining
both summands and their
original ratio.

If $v(a)<v(b)$, that
ratio $b/a$ belongs to
$\mathfrak m_A$. Thus
$1+b/a$ is a unit:
if it were in the maximal
ideal, subtracting $b/a$
would put $1$ there.
Consequently $a+b\neq0$
and $v(a+b)=v(a)$.
The reversed strict inequality
gives the symmetric equality.
This entire proof also
applies to nonzero $a,b\in K$
outside $A$: the comparison
still puts the same ratio
in $A$. The original value
map on $K^*$ therefore
satisfies the same inequality,
with equality whenever the
two original values differ.

For a finite original list
$a_1,\ldots,a_n\in K$, keep
all its terms, including
zeros. If at least one
term is nonzero, choose
$a_j$ of least value among
the nonzero terms. Every
ratio $a_i/a_j$ belongs
to $A$, with a zero
term giving the ratio $0$.
Hence if the original sum
$a=\sum_i a_i$ is nonzero,
then $a/a_j=\sum_i a_i/a_j$
is a nonzero element of
$A$, proving $v(a)\geq v(a_j)$.
If exactly one term has
the least value, every
other ratio is in
$\mathfrak m_A$, including
the zeros. The full sum
of ratios is $1$ plus
an element of that ideal,
hence a unit. In that
case the sum is nonzero
and has exactly the least
original value. For an
all-zero or empty list,
the sum is $0$ and no
value in $\Gamma$ is
assigned to it.
\end{proof}

\begin{lemma}
\label{editorial-source-lemma-characterize-valuation-ring}
Let $A$ be a ring. The following are equivalent
\begin{enumerate}
\item $A$ is a valuation ring,
\item $A$ is a local domain and every finitely generated
ideal of $A$ is principal.
\end{enumerate}
\end{lemma}

\begin{proof}
Suppose first that $A$ is
a valuation ring. Retain
the complete finite generator
list $f_1,\ldots,f_n$.
If the list is empty
or all generators are zero,
its ideal is $(0)$ and
is principal. Otherwise choose
an original nonzero $f_i$
whose value is least among
the nonzero generators, and
put $h=f_i$. For each
nonzero $f_j$, the comparison
$v(f_j)\geq v(h)$ gives
the original ratio $c_j=f_j/h\in A$.
For $f_j=0$, set $c_j=0$.
Then the exact equalities
$f_j=hc_j$ for every $j$
give $(f_1,\ldots,f_n)\subseteq(h)$,
while $h$ itself is in
the original list, giving
the reverse inclusion.
No value is taken at
a zero generator.

Conversely, assume the original
$A$ is a local domain
with every finitely generated
ideal principal. An original
fraction $f/g$ with $g\neq0$
belongs to $A$ if $f=0$.
For $f,g\neq0$, write
$(f,g)=(h)$. Then $h\neq0$,
and there are original
$a,b,c,d\in A$ with
$$
f=ah,\qquad g=bh,\qquad h=cf+dg.
$$
Substitution gives
$h=(ca+db)h$. The domain
property and $h\neq0$ permit
cancellation, giving $ca+db=1$.
Both $a,b$ cannot belong
to the maximal ideal,
so at least one is a
unit. If $a$ is a
unit, the original fraction
$g/f=b/a$ belongs to $A$;
if $b$ is a unit,
$f/g=a/b$ belongs to $A$.
This proves the nonzero
fraction alternative in the
actual $\operatorname{Frac}(A)$.
Lemma
\ref{editorial-source-lemma-x-or-x-inverse-valuation-ring}
therefore makes $A$ a
valuation ring.

The converse proof uses
only two-generated ideals.
Thus in (2) it suffices
to require that every
two-generated ideal be principal;
the forward proof then
recovers the assertion for
every original finite list.
The local hypothesis cannot
be discarded: every ideal
of $\mathbf Z$ is principal,
since for a nonzero ideal
its least positive element
divides all its elements
by division with remainder,
and the zero ideal is
$(0)$. But neither $2/3$
nor $3/2$ is an integer,
so the exact fraction
criterion shows that $\mathbf Z$
is not a valuation ring.
\end{proof}

\begin{lemma}
\label{editorial-source-lemma-valuation-valuation-ring}
Let $(\Gamma, \geq)$ be a totally ordered abelian group.
Let $K$ be a field. Let $v : K^* \to \Gamma$ be a homomorphism
of abelian groups such that $v(a + b) \geq \min(v(a), v(b))$ for
$a, b \in K$ with $a, b, a + b$ not zero. Then
$$
A =
\{
x \in K \mid x = 0 \text{ or } v(x) \geq 0
\}
$$
is a valuation ring with value group $\Im(v) \subset \Gamma$,
with maximal ideal
$$
\mathfrak m =
\{
x \in K \mid x = 0 \text{ or } v(x) > 0
\}
$$
and with group of units
$$
A^* =
\{
x \in K^* \mid v(x) = 0
\}.
$$
\end{lemma}

\begin{proof}
Keep the original ordered group
$\Gamma$ and homomorphism
$v:K^*\to\Gamma$, without
assuming it is surjective.
The homomorphism identities give
$v(1)=0$ and $v(x^{-1})=-v(x)$.
Also $2v(-1)=v(1)=0$.
A totally ordered abelian group
has no nonzero element killed
by a positive integer, by
the strict-order argument in
Lemma \ref{editorial-source-lemma-valuation-group}.
Hence $v(-1)=0$, and
$v(-x)=v(x)$ for nonzero $x$.

The original set $A$ contains
$0,1$ and is closed under
additive inverses. It is
closed under multiplication:
products involving zero are
zero, and otherwise the
sum of two nonnegative
original values is nonnegative.
For addition, a zero summand
causes no problem. If both
summands and their sum are
nonzero, the given inequality
puts the sum in $A$;
if their sum is zero,
it belongs to $A$ by
definition. Thus $A$ is
a unital subring of $K$.
For every nonzero $x\in K$,
totality gives $v(x)\geq0$
or $v(x)<0$. In the
second case $v(x^{-1})>0$,
so $x^{-1}\in A$.
The criterion of Lemma
\ref{editorial-source-lemma-x-or-x-inverse-valuation-ring}
therefore proves that $A$
is a valuation ring whose
fraction field is exactly $K$.

A nonzero original $x$
is a unit of $A$ exactly
when both $v(x)$ and
$-v(x)$ are nonnegative,
which by antisymmetry is
equivalent to $v(x)=0$.
Thus its nonunits are
precisely $0$ and the
elements of positive value,
giving the asserted maximal
ideal. Directly, that set
is a proper ideal: addition
uses the minimum inequality
unless the sum is zero,
and multiplying a positive
value by an element of
$A$ adds a nonnegative
value. It excludes $1$,
and every element outside
it has the inverse just
described. This also proves
maximality explicitly.

Let $q:K^*\to\Delta=K^*/A^*$
be the original associated
value-group quotient. The map
$$
\overline v:\Delta\longrightarrow\operatorname{Im}(v),
\qquad[x]\longmapsto v(x)
$$
is well-defined because
$A^*=\ker(v)$, injective
by the same kernel equality,
and surjective by the
definition of the image.
Its inverse sends a value
$\gamma=v(x)$ to $[x]$;
different original choices
$x,y$ have $v(x/y)=0$
and differ by a unit of
$A$, proving independence.
Moreover
$$
[x]\geq[y]\quad\Longleftrightarrow\quad
x/y\in A\quad\Longleftrightarrow\quad
v(x)-v(y)\geq0.
$$
Thus it is an order
isomorphism with exactly
the asserted subgroup of
the original $\Gamma$.
If $i:\operatorname{Im}(v)\hookrightarrow\Gamma$
is its original inclusion,
then $i\overline vq=v$;
the ambient codomain and
the original value map
have both been retained.

Finally two such original
maps $v:K^*\to\Gamma$ and
$w:K^*\to\Lambda$ define
the same valuation ring
exactly when there is an
order isomorphism
$\phi:\operatorname{Im}(v)\to\operatorname{Im}(w)$
with $\phi(v(x))=w(x)$
for every original $x$.
For a common ring, equality
of two $v$-values means
their ratio is a unit
of that ring, and hence
also has $w$-value zero.
This proves the formula
well-defined; the reverse
argument gives its inverse.
The same ratio criterion
proves that it preserves
and reflects order. Conversely
such an isomorphism preserves
the nonnegative cone, so
the two original sets
$\{0\}\cup v^{-1}(\Gamma_{\geq0})$
and $\{0\}\cup w^{-1}(\Lambda_{\geq0})$
are equal. Uniqueness follows
because every element of
the image is an original
value $v(x)$. No isomorphism
between unused parts of the
ambient groups is asserted.
\end{proof}

\noindent
Let $(\Gamma, \geq)$ be a totally ordered abelian group.
An {\it ideal of $\Gamma$} is a subset $I \subset \Gamma$ such
that all elements of $I$ are $\geq 0$ and $\gamma \in I$,
$\gamma' \geq \gamma$ implies $\gamma' \in I$. We say that such
an ideal is {\it prime} if $0 \not \in I$ and if
$\gamma + \gamma' \in I, \gamma, \gamma' \geq 0
\Rightarrow \gamma \in I \text{ or } \gamma' \in I$.

\begin{lemma}
\label{editorial-source-lemma-ideals-valuation-ring}
Let $A$ be a valuation ring.
Ideals in $A$ correspond $1 - 1$ with ideals of $\Gamma$.
This bijection is inclusion preserving, and maps prime
ideals to prime ideals.
\end{lemma}

\begin{proof}
Retain $K=\operatorname{Frac}(A)$, the original ordered quotient
$\Gamma=K^*/A^*$, and its quotient map $v:K^*\to\Gamma$.
For an ideal $J\subset A$ and an ideal $U\subset\Gamma$, define
$$
V(J)=\{v(x):x\in J,\ x\neq0\},\qquad
J(U)=\{0\}\cup\{x\in K^*:v(x)\in U\}.
$$
Every element of $V(J)$ is nonnegative. If $v(x)=\gamma\in V(J)$
and $\delta\geq\gamma$, choose $y\in K^*$ with $v(y)=\delta$.
Then $v(y/x)=\delta-\gamma\geq0$, so $y/x\in A$ and
$y=(y/x)x\in J$. Thus $V(J)$ is upward closed.
The same calculation shows that every representative of a value in
$V(J)$ lies in $J$; when the values are equal, $y/x$ is an actual unit.

Since $U\subset\Gamma_{\geq0}$, the set $J(U)$ lies in $A$.
It contains zero. If $x\neq0$ belongs to it, then
$v(-x)=v(x)\in U$. For nonzero $x,y\in J(U)$ with $x+y\neq0$,
the valuation inequality gives
$v(x+y)\geq\min\{v(x),v(y)\}\in U$, hence $x+y\in J(U)$.
If a summand or the sum is zero the closure follows directly from
the definition. For $a\in A$ and $x\in J(U)$, a zero factor gives
$ax=0$; otherwise $v(ax)=v(a)+v(x)\geq v(x)$, so $ax\in J(U)$.
This proves that $J(U)$ is an ideal.
Surjectivity of $v$ and the representative calculation give
$V(J(U))=U$ and $J(V(J))=J$. Both maps preserve inclusion.
In particular $J(\varnothing)=(0)$ and
$J(\Gamma_{\geq0})=A$; no value is assigned to the zero element.

The ideal $J$ is proper exactly when $0\notin V(J)$, since an
element of value zero is a unit. If $J$ is prime and
$\gamma,\delta\geq0$ with $\gamma+\delta\in V(J)$, choose
$x,y\in A\setminus\{0\}$ with these values. The representative
calculation puts $xy$ in $J$, so $x\in J$ or $y\in J$.
Consequently $\gamma\in V(J)$ or $\delta\in V(J)$.
Conversely, suppose $U$ is prime and $xy\in J(U)$ for $x,y\in A$.
If $xy=0$, the domain property gives a zero factor, which belongs
to $J(U)$. Otherwise the defining condition on
$v(x)+v(y)\in U$ puts one factor in $J(U)$.
Thus the correspondence preserves and reflects primality, including
the empty group ideal and the zero ring ideal.

The correspondence also describes the ideal operations exactly.
For every family $(J_\lambda)_{\lambda\in\Lambda}$, put
$U_\lambda=V(J_\lambda)$. Then
$$
V\left(\bigcap_{\lambda\in\Lambda}J_\lambda\right)
=\bigcap_{\lambda\in\Lambda}U_\lambda,
\qquad
V\left(\sum_{\lambda\in\Lambda}J_\lambda\right)
=\bigcup_{\lambda\in\Lambda}U_\lambda.
$$
The first equality follows by membership of each nonzero element.
For the second, a nonzero member of the sum is a finite sum of
members of the original ideals. Discarding no term from that
equality, list its nonzero terms. Their list is nonempty; its least
value belongs to one $U_\lambda$. The finite-sum inequality and
upward closure put the value of the sum in that same $U_\lambda$.
The reverse inclusion follows from each $J_\lambda$ being contained
in the sum. For an empty family the intersection is $A$ and its
value set is $\Gamma_{\geq0}$, while the sum is $(0)$ and its
value set is empty; the intersections on the right are taken
inside $\Gamma_{\geq0}$.
For two ideals with value sets $U,W$, one has
$$
V(J(U)J(W))=U+W
=\{\gamma+\delta:\gamma\in U,\ \delta\in W\},
$$
and
$$
V(\sqrt{J(U)})
=\{\gamma\in\Gamma_{\geq0}:n\gamma\in U
\text{ for some integer }n\geq1\}.
$$
Indeed $U+W$ is upward closed: if
$\eta\geq\gamma+\delta$, write
$\eta=\gamma+(\delta+\eta-\gamma-\delta)$ and use
$\delta+\eta-\gamma-\delta\geq\delta$.
Every value in $U+W$ is the value of a product of representatives
from the original ideals. A nonzero element of their product ideal
is a finite sum of such products, and the least-value argument
just used gives the reverse inclusion. If either ideal is zero,
both sides are empty. Finally, for nonzero $x\in A$, the equality
$v(x^n)=n v(x)$ proves the radical formula in both directions;
the original zero element lies in the radical separately.
\end{proof}

\begin{lemma}
\label{editorial-source-lemma-valuation-ring-Noetherian-discrete}
A valuation ring is Noetherian if and only if it is
a discrete valuation ring or a field.
\end{lemma}

\begin{proof}
Let $K=\operatorname{Frac}(A)$ and retain the original value map
$v:K^*\to\Gamma=K^*/A^*$. If $A$ is a field, its two ideals
are $(0)$ and $A$, so it is Noetherian.
If $A$ is a discrete valuation ring, let $\gamma_0>0$ generate
its original ordered group, and let
$\phi:\Gamma\to\mathbf Z$ be the order isomorphism
$\phi(n\gamma_0)=n$. Choose $\pi\in A$ with
$v(\pi)=\gamma_0$. The nonzero ideals, and only the nonzero
ideals, are
$$
I_n=\{0\}\cup\{x\in K^*:v(x)\geq n\gamma_0\}
=(\pi^n),\qquad n\geq0.
$$
To verify this directly, the values of a nonzero ideal under
$\phi$ form a nonempty upward closed subset of
$\mathbf Z_{\geq0}$, with a least element $n$.
For an original element $a$ of that value,
$u=a/\pi^n$ has value zero, so $u\in A^*$ and
$a=u\pi^n$. Every original nonzero $b$ of the ideal satisfies
$b/\pi^n\in A$, while $0=\pi^n0$. These equalities prove
that the ideal is $(\pi^n)$ and retain its original unit factor.
The remaining ideal $(0)$ is generated by zero. Thus every ideal
is finitely generated and $A$ is Noetherian. The map $v$ still
takes values in $\Gamma$; the integer-valued comparison is
$\phi\circ v$, with $v=\phi^{-1}\circ(\phi\circ v)$.

Conversely, suppose $A$ is Noetherian and is not a field.
Ascending chains of its ideals stabilize: their union is an ideal,
and finitely many generators of that union all occur in a single
member of the chain. By Lemma \ref{editorial-source-lemma-ideals-valuation-ring},
the same holds for the upward closed subsets of
$\Gamma_{\geq0}$. There is a positive element of $\Gamma$.
If there were no least positive element, starting from one such
element one could choose
$\gamma_1>\gamma_2>\gamma_3>\cdots>0$.
The corresponding ideals
$\gamma_i+\Gamma_{\geq0}$ would form a strictly ascending
chain: $\gamma_{i+1}$ belongs to its own set but not to the
preceding one. This is impossible. Denote the least positive
element by $\gamma_0$.

For any original $\gamma>0$, suppose that
$\gamma-n\gamma_0\geq0$ for every integer $n\geq0$.
These elements strictly decrease, and their upward closed sets
again give a strictly ascending chain of ideals, a contradiction.
Consequently there is a least integer $N\geq1$ such that
$\gamma-N\gamma_0<0$. Put
$r=\gamma-(N-1)\gamma_0$. Then
$0\leq r<\gamma_0$, so the least-positive property forces
$r=0$. Thus $\gamma=(N-1)\gamma_0$.
Zero is $0\gamma_0$, and applying the argument to $-\gamma$
expresses every negative element as a negative multiple of
$\gamma_0$. Distinct integer multiples are distinct, because
a positive multiple of $\gamma_0$ is positive. Therefore
$n\mapsto n\gamma_0$ is an order isomorphism
$\mathbf Z\to\Gamma$, and $A$ is a discrete valuation ring.
The argument uses the ascending chain condition twice; merely
having a least positive group element is not the assertion used
to conclude that all positive elements are its multiples.
\end{proof}























